\documentclass[12pt]{article}
\usepackage{amsmath,amssymb}
\usepackage{geometry}
\usepackage{enumitem}
\geometry{margin=0.8in}
\setlength{\parindent}{0pt}

\begin{document}

\begin{center}
  {\Large\textbf{Trigonometric Identity Proofs}}\\[3pt]
  {\large Precalculus Practice --- 45 Minutes}
\end{center}

Name: \rule{2.7in}{0.4pt}\hfill Date: \rule{1.5in}{0.4pt}

\section*{Instructions}
Prove each identity by transforming the left-hand side into the right-hand side.
Show a chain of equal expressions and name the identity or algebraic rule used
at each step. Do not manipulate both sides as an equation or verify with
numerical examples. Use exact expressions; no calculator. All angles are in
radians. Assume variables take values for which both sides are defined; observe
any additional interval restriction.

This set covers cofunction, sum/difference, double-angle, half-angle, and
even/odd identities. Basic reciprocal, quotient, and Pythagorean identities
may also be used.

\section*{Prove each identity}
\begin{enumerate}[leftmargin=*,itemsep=0pt]
  \item
  \[
    \frac{\cos\left(\frac{\pi}{2}-x\right)}{\cos(-x)}=\tan x
  \]
  \vspace{\stretch{1}}

  \item
  \[
    \tan(-x)\csc(-x)=\sec x
  \]
  \vspace{\stretch{1}}

  \item
  \[
    \sin(x+y)+\sin(x-y)=2\sin x\cos y
  \]
  \vspace{\stretch{1}}

  \item
  \[
    \frac{1-\cos(2x)}{\sin(2x)}=\tan x
  \]
  \vspace*{\stretch{1}}
\end{enumerate}

\newpage
\begin{center}
  {\large\textbf{Trigonometric Identity Proofs --- Continued}}
\end{center}

\begin{enumerate}[leftmargin=*,itemsep=0pt,resume]
  \item
  \[
    \frac{\sec\left(\frac{\pi}{2}-x\right)}{\cot x}=\sec x
  \]
  \vspace{\stretch{1}}

  \item
  \[
    \frac{\tan(x+y)-\tan x}{\tan y}
    =\frac{1+\tan^2x}{1-\tan x\tan y}
  \]
  \vspace{\stretch{1}}

  \item
  \[
    \frac{\sin(-x)+\cos\left(\frac{\pi}{2}+x\right)}{\cos(-x)}
    =-2\tan x
  \]
  \vspace{\stretch{1}}

  \item
  \[
    \frac{\cos(2x)}{\cos x-\sin x}=\cos x+\sin x
  \]
  \vspace*{\stretch{1}}
\end{enumerate}

\newpage
\begin{center}
  {\large\textbf{Trigonometric Identity Proofs --- Continued}}
\end{center}

\begin{enumerate}[leftmargin=*,itemsep=0pt,resume]
  \item
  \[
    \frac{\sin^2\left(\frac{x}{2}\right)}{1+\cos x}
    =\frac12\tan^2\left(\frac{x}{2}\right)
  \]
  \vspace{\stretch{1}}

  \item For $0<x<\pi$,
  \[
    \sqrt{\frac{1-\cos x}{1+\cos x}}=\tan\left(\frac{x}{2}\right).
  \]
  \vspace{\stretch{1}}

  \item
  \[
    \frac{\sin(3x)-\sin x}{\cos(3x)+\cos x}=\tan x
  \]
  \vspace{\stretch{1}}

  \item
  \[
    \frac{\sin(x+y)}{\sin(x-y)}
    =\frac{\tan x+\tan y}{\tan x-\tan y}
  \]
  \vspace*{\stretch{1}}
\end{enumerate}

\end{document}
